\icode{E-EMI-CA-L3-261004-002}
如图所示，两根光滑平行导轨由半径为\(r\)的四分之一绝缘圆弧部分与不计电阻的水平金属部分组成，圆弧部分与水平部分平滑连接并固定在水平地面上，导轨间距为\(L\)，水平导轨部分存在竖直向下的匀强磁场，磁感应强度大小为\(B\)。一质量为\(2m\)、电阻为\(R\)的金属棒\(ab\)静置于水平导轨与圆轨道的连接处，另一质量为\(m\)、电阻为\(2R\)的金属棒\(cd\)从圆弧最高处静止释放，下滑至水平轨道后与\(ab\)棒发生弹性正碰。若金属棒\(cd\)与金属棒\(ab\)始终垂直于金属导轨并接触良好，水平轨道足够长，重力加速度为\(g\)，求：
\begin{enumerate}
	\item 金属棒\(cd\)滑到轨道底部与\(ab\)棒碰撞前，导轨对金属棒\(cd\)的支持力；
	\item 金属棒\(cd\)进入磁场后运动的最大速度及整个过程中\(ab\)棒产生的热量；
	\item 若在金属棒\(cd\)达到最大速度时两棒间距为\(x_0\)，给金属棒\(ab\)施加一水平向右的恒力\(F\)，当两棒运动稳定后两棒间距\(\Delta x\)与时间\(t\)的关系。
\end{enumerate}
\begin{tikzpicture}[font=\scriptsize]
	\draw (-1,1) arc(180:270:1)--++(4,0);
	\draw ($(-1,1)+(40:1)$) arc(180:270:1)--++(4,0);
	\draw[very thick,red] ($(0,0)+(40:-.2)$)node[below left]{\(a\)}--++(40:1.4)node[above right]{\(b\)} ($(-1,1)+(40:-.2)$)node[below left]{\(c\)}--++(40:1.4)node[above right]{\(d\)};
	\foreach \x in {1,1.8,...,4.7}
		\draw[->,blue] (\x,1.2)--++(0,-.9);
\end{tikzpicture}
\begin{answer}
	（1）\(3mg\)（2）\(\frac{2}{81}mgr\)（3）\(\Delta x = x_0 + \frac{RF}{B^2L^2}\cdot t\)
\end{answer}